%%%%%%%% Cost-Sensitive Conformal Decision-Making Under Label-Shift Uncertainty %%%%%%%%%

\documentclass{article}

\usepackage{microtype}
\usepackage{graphicx}
\usepackage{subcaption}
\usepackage{booktabs}
\usepackage{hyperref}

\newcommand{\theHalgorithm}{\arabic{algorithm}}

\usepackage[accepted]{icml2026}

\renewcommand{\ttdefault}{cmtt}

\usepackage{amsmath}
\usepackage{amssymb}
\usepackage{mathtools}
\usepackage{amsthm}
\usepackage{algorithm}
\usepackage{algorithmic}

\usepackage[capitalize,noabbrev]{cleveref}

%%%%%%%%%%%%%%%%%%%%%%%%%%%%%%%%
% THEOREMS
%%%%%%%%%%%%%%%%%%%%%%%%%%%%%%%%
\theoremstyle{plain}
\newtheorem{theorem}{Theorem}[section]
\newtheorem{proposition}[theorem]{Proposition}
\newtheorem{lemma}[theorem]{Lemma}
\newtheorem{corollary}[theorem]{Corollary}
\theoremstyle{definition}
\newtheorem{definition}[theorem]{Definition}
\newtheorem{assumption}[theorem]{Assumption}
\theoremstyle{remark}
\newtheorem{remark}[theorem]{Remark}

%%%%%%%%%%%%%%%%%%%%%%%%%%%%%%%%
% NOTATION — consistent conventions:
%   Capitals (X, Y, R): random variables / expected quantities
%   Calligraphic (\mathcal{X}): sets / spaces
%   Greek (\alpha, \lambda, \rho): parameters / thresholds
%   Lowercase (c, w, n): constants / realised values
%   Subscripts: class labels (y), distribution (tr/dep)
%   Superscripts: distribution (tr/dep) on priors
%   i in {1,...,N}: sample index
%%%%%%%%%%%%%%%%%%%%%%%%%%%%%%%%
\newcommand{\loss}{\ell}                         % per-sample loss
\newcommand{\Risk}{R}                            % expected risk
\newcommand{\Rhat}{\hat{R}}                      % empirical risk
\newcommand{\lamhat}{\hat{\lambda}}              % selected threshold
\newcommand{\pii}[2]{\pi_{#1}^{\mathrm{#2}}}    % class prior: \pii{y}{tr}
\newcommand{\rhostar}{\rho^{*}}                  % true shift ratio
\newcommand{\rhohat}{\hat{\rho}}                 % assumed shift ratio
\newcommand{\Bpw}{B_{\mathrm{pw}}}              % prevalence-weighted bound
\newcommand{\Btv}{B_{\mathrm{tv}}}              % total variation bound
\newcommand{\asafe}{\alpha_{\mathrm{safe}}}      % adjusted calibration level
\newcommand{\cfn}{c_{\mathrm{fn}}}               % false negative cost
\newcommand{\cfp}{c_{\mathrm{fp}}}               % false positive cost

\icmltitlerunning{When Should an LLM Say ``I Don't Know''?}

\begin{document}

\twocolumn[
  \icmltitle{When Should an LLM Say ``I Don't Know''?:\\
    Conformal Risk Control for Safety-Critical\\
    LLM Deployment Under Shift}

  \icmlsetsymbol{equal}{*}

  \begin{icmlauthorlist}
    \icmlauthor{Baljinnyam Dayan}{ic}
  \end{icmlauthorlist}

  \icmlaffiliation{ic}{Imperial College London, CID: 06059151}
  \icmlcorrespondingauthor{Baljinnyam Dayan}{bd1125@ic.ac.uk}

  \icmlkeywords{LLM safety, Hallucination, Abstention, Conformal prediction, Label shift, Risk control, Agentic AI, GAIA benchmark}

  \vskip 0.3in
]

\printAffiliationsAndNotice{}

% ======================================================================
\begin{abstract}
  Large language models are increasingly used in settings where a confident wrong answer is worse than no answer at all: medical assistants, safety filters, and tool-using copilots acting on behalf of users.
  Recent work argues that hallucinations persist in part because current training and evaluation reward guessing over acknowledging uncertainty.
  We study the complementary \emph{deployment-time} problem: given a scored LLM output, when should the system answer, retry, or say ``I don't know'' / hand off under distribution shift?
  We extend Conformal Risk Control (CRC) to cost-sensitive decision-making under uncertain label shift with four contributions:
  (1)~a prevalence-weighted excess risk bound that is provably tighter than the standard total variation bound (27--39\% improvement);
  (2)~an \emph{uncertainty-aware recalibration} procedure with a \emph{three-way decision rule} that maintains cost-weighted guarantees over a practitioner-specified or data-driven shift interval, with abstention emerging from endpoint threshold disagreement and a formal conditional risk guarantee on the decided subset;
  (3)~empirical evaluation on two LLM-facing deployment settings---toxicity moderation (CivilComments, $\cfn/\cfp = 5$) and agentic QA on GAIA---with Folktables as non-LLM support, showing that recalibration eliminates risk violations under moderate shift, while Bayesian three-way decisions convert uncertainty into automatic escalation;
  (4)~an \emph{agentic deployment layer} on GAIA that extends CRC from classifier scores to LLM-agent badness scores, reducing autonomous risk on a controlled difficulty-shift demo and abstaining safely when the underlying agent lacks sufficient capability on the full benchmark.
  The resulting perspective is practical rather than magical: we do not solve hallucination at training time, but we provide a principled deployment layer that helps an LLM refuse unsafe answers.
\end{abstract}

% ======================================================================
% BEAT 1: Motivation — why risk guarantees matter
% BEAT 2: Issue — label shift + cost asymmetry + unknown shift = unsolved
% BEAT 3: Contribution — what we do about it
% ======================================================================
\section{Introduction}
\label{sec:intro}

% --- Beat 1: Motivation ---
Large language models are increasingly deployed in settings where a confident wrong answer can be dangerous: medical information assistants, safety filters, and tool-using copilots that act on behalf of users.
Yet modern systems are usually optimized to \emph{produce an answer}, not to know when they should abstain.
Recent work argues that hallucinations persist in part because training and evaluation reward guessing over acknowledging uncertainty \citep{kalai2025hallucinate}.
Conformal Risk Control \citep[CRC;][]{angelopoulos2024conformal} offers finite-sample risk guarantees for scored decisions, but assumes exchangeable calibration and deployment data.

% --- Beat 2: Issue and why it remains unsolved ---
The deployment distribution shifts constantly.
A toxicity filter calibrated on general comments ($6\%$ toxic) faces identity-targeting discussions ($29\%$ toxic); an LLM assistant calibrated on easy or text-only tasks may face harder tasks, attachments, or agentic workflows at deployment.
These are settings where the rate of bad outcomes changes, and where the central operational question is not only \emph{what} answer to produce but \emph{whether to answer at all}.
Under classical label shift, the marginal $P(Y)$ changes while the class-conditional $P(X \mid Y)$ does not; in agentic settings, the closest analogue is a \emph{failure-prevalence shift} in the rate of wrong first-pass answers.
Prior work addresses label shift through importance weighting \citep{tibshirani2019conformal, podkopaev2021distribution} or PAC prediction sets \citep{fisch2022conformal}, but two practical gaps remain.
First, existing guarantees target symmetric losses (miscoverage, set size) rather than the \emph{cost-weighted} risk relevant to safety-critical LLM deployment, where a dangerous hallucination can cost much more than an unnecessary abstention.
Second, importance weighting requires knowing---or accurately estimating---the deployment prevalence.
When the deployment distribution is uncertain, the practitioner faces a fundamental question: \textbf{when should an LLM answer, when should it retry, and when should it say ``I don't know'' / hand off?}

% --- Beat 3: Our contribution ---
We address all three gaps.
We make four contributions:
\begin{enumerate}
  \item A \textbf{prevalence-weighted excess risk bound} that exploits cost--prevalence structure to tighten the standard total variation bound, with improvement governed by class imbalance and cost ratio (\cref{sec:bound}; 27--39\% in our settings, formally characterised in \cref{cor:regime}).
  \item \textbf{Uncertainty-aware recalibration with three-way decisions}: a procedure that computes the tightest calibration level guaranteeing cost-weighted risk control for \emph{all} shift ratios in a practitioner-specified or data-driven interval $[\rho_{\mathrm{lo}}, \rho_{\mathrm{hi}}]$, with a three-way decision rule (``classify / defer / escalate'') whose deferral emerges from endpoint threshold disagreement---no deferral cost hyperparameter---and a formal conditional risk guarantee on the decided subset (\cref{thm:threeway,sec:b4}).
  \item An \textbf{empirical study} centred on safety-critical LLM deployment: toxicity moderation on CivilComments ($\cfn/\cfp = 5$) and agentic QA on GAIA, with Folktables as non-LLM support.  Across these settings, recalibration improves safety margins under moderate shift, three-way decisions provide automatic escalation, and conservative recalibration can distribute burden unevenly across subgroups (\cref{sec:results}).
  \item An \textbf{agentic deployment layer} on the GAIA benchmark \citep{mialon2023gaia}, applying CRC to LLM-agent \emph{badness scores} (LLM-as-judge estimates of $P(\text{answer is wrong})$) under task-difficulty shift.  In a controlled difficulty-shift demo, CRC-based gating reduces autonomous risk from $37.5\%$ to $20.0$%; on the full benchmark, it safely abstains when the underlying agent lacks sufficient capability (\cref{sec:gaia_results}).
\end{enumerate}

% ======================================================================
\section{Related Work}
\label{sec:related}

% --- Paragraph 1: Conformal prediction under distribution shift ---
\textbf{Conformal methods under distribution shift.}
Several approaches extend conformal guarantees beyond exchangeability.
\citet{tibshirani2019conformal} introduce weighted conformal prediction for covariate shift using likelihood ratios.
\citet{podkopaev2021distribution} provide coverage guarantees under label shift via importance-weighted quantiles.
\citet{barber2023conformal} develop a general framework for conformal prediction beyond exchangeability, covering both covariate and label shift.
More recently, \citet{angelopoulos2026nonmonotonic} extends conformal risk control to non-monotonic, multi-dimensional losses, and \citet{hultberg2026anytime} provide anytime-valid CRC guarantees over growing calibration sets.
However, these methods target coverage or general risk functionals, not \emph{cost-weighted} decision risk under uncertain shift magnitude.

% --- Paragraph 2: Label shift estimation ---
\textbf{Label shift estimation.}
Detecting and correcting for label shift requires estimating the deployment class distribution.
\citet{lipton2018detecting} propose Black Box Shift Estimation (BBSE), which uses the classifier's confusion matrix and unlabelled deployment predictions.
\citet{alexandari2020maximum} improve upon BBSE with maximum-likelihood calibration.
\citet{saerens2002adjusting} give an EM procedure for adjusting classifier outputs to new priors.
These methods provide point estimates of the shift, but do not quantify the downstream risk when the estimate is wrong---a gap our bound (\cref{sec:bound}) and recalibration procedure (\cref{sec:b4}) address.

% --- Paragraph 3: PAC prediction sets ---
\textbf{PAC prediction sets under label shift.}
\citet{fisch2022conformal} provide PAC-style guarantees for prediction sets under label shift, combining conformal calibration with worst-case bounds over a family of distributions.
\citet{xu2025selective} propose Selective Conformal Risk Control, integrating conformal prediction with selective classification via a two-stage confidence-then-calibrate pipeline with PAC-style guarantees.
We extend the \citet{fisch2022conformal} framework from prediction-set coverage to cost-sensitive threshold decisions with uncertainty-aware recalibration.

% --- Paragraph 4: Fairness in conformal prediction ---
\textbf{Fairness in conformal prediction.}
Recent work studies equalized coverage across demographic groups \citep{romano2020malice}, and \citet{vadlamani2025conformal} propose a generic framework adapting group fairness notions (demographic parity, equalized odds) to conformal prediction.
Our subgroup analysis (\cref{sec:results}) complements this line by showing that \emph{aggregate} cost-weighted risk control can mask unequal false-negative burden across groups---a distinct failure mode from unequal coverage.

% --- Paragraph 5: Cost-sensitive conformal methods and deferral ---
\textbf{Cost-sensitive conformal methods and deferral.}
Concurrent work addresses cost-sensitivity and deferral within conformal frameworks.
\citet{jia2026costsensitive} propose cost-sensitive conformal \emph{training} via rank weighting, optimising prediction-set size end-to-end during model training rather than post-hoc calibration.
Most closely related, \citet{kwon2026conformal} combine conformal prediction with cost-aware deferral for clinical triage under distribution shift, using importance weighting for shift robustness and a cost-based deferral rule that reduces error by ${\sim}50\%$ on retained cases.
Our work differs in three respects: (i)~we derive a \emph{tighter} prevalence-weighted excess risk bound (\cref{thm:bound}) rather than relying on standard importance weighting; (ii)~our three-way deferral emerges from shift-interval disagreement with \emph{zero free parameters}, whereas their deferral requires a cost hyperparameter; and (iii)~we provide a formal conditional risk guarantee on the decided subset (\cref{thm:threeway}), which their framework does not.

% --- Paragraph 6: Conformal methods for language models ---
\textbf{Hallucination, abstention, and conformal methods for language models.}
Recent work argues that LLM hallucinations persist in part because training and evaluation reward guessing over acknowledging uncertainty \citep{kalai2025hallucinate}.
Our work is complementary: rather than modifying pretraining or benchmark design, we study the \emph{deployment-time} question of when a scored LLM should answer, retry, or abstain.
Recent conformal work on language models includes guaranteed generation quality \citep{quach2024conformal}, cost-efficient cascades \citep{mohri2024language}, and conformal methods for safe generation or QA \citep{kumar2023conformal}.
These methods focus on generation quality or answer selection; we target \emph{safety-critical abstention} under shift, where the core requirement is that a deployed LLM should know when not to answer.

% ======================================================================
\section{Background}
\label{sec:background}

We introduce notation and review the two building blocks our method extends: Conformal Risk Control and importance-weighted calibration under label shift.
We summarise the key symbols in \cref{tab:notation}.

\begin{table}[t]
  \caption{Notation summary.}
  \label{tab:notation}
  \begin{center}
    \begin{small}
      \begin{tabular}{cl}
        \toprule
        Symbol                             & Meaning                                        \\
        \midrule
        $X, Y$                             & Feature and label random variables             \\
        $f\!: \mathcal{X} \to [0,1]$       & Scoring function                               \\
        $\lambda \in [0,1]$                & Decision threshold                             \\
        $\lamhat$                          & Threshold selected by CRC                      \\
        $(x_i, y_i)_{i=1}^{N}$             & Calibration data ($N$ samples)                 \\
        $\cfn,\, \cfp$                     & False-negative / false-positive cost           \\
        $\alpha$                           & Target risk level                              \\
        $\pii{y}{tr},\, \pii{y}{dep}$      & Class $y$ prior (train / deploy)               \\
        $\rhostar$                         & True label-shift odds ratio                    \\
        $\rhohat$                          & Assumed shift ratio                            \\
        $w_y$                              & Importance weight $\pii{y}{dep} / \pii{y}{tr}$ \\
        $\Bpw,\, \Btv$                     & Bound constants (ours / TV)                    \\
        $\asafe$                           & Adjusted calibration level (recalibration)     \\
        $\lambda_{\min},\, \lambda_{\max}$ & Three-way defer zone boundaries                \\
        \bottomrule
      \end{tabular}
    \end{small}
  \end{center}
  \vskip -0.1in
\end{table}

\subsection{Conformal Risk Control}
\label{sec:crc}

Let $f\!: \mathcal{X} \to [0,1]$ be a fixed scoring function and $\lambda \in [0,1]$ a decision threshold, so the prediction is
\begin{equation}
  \hat{Y} = \mathbf{1}[f(X) \geq \lambda].
\end{equation}
Given calibration data $(x_i, y_i)_{i=1}^{N}$ drawn exchangeably from~$P$, the cost-weighted loss is:
\begin{align}
  \loss(\lambda; x, y)
   & = \cfn \cdot \mathbf{1}[f(x) < \lambda,\; y = 1] \notag   \\
   & \quad + \cfp \cdot \mathbf{1}[f(x) \geq \lambda,\; y = 0]
  \label{eq:loss}
\end{align}
where $\cfn > \cfp > 0$ reflects asymmetric error costs.
CRC \citep{angelopoulos2024conformal} selects the smallest $\lamhat$ satisfying:
\begin{equation}
  \Rhat(\lamhat) + \frac{b - \Rhat(\lamhat)}{N+1} \leq \alpha
  \label{eq:crc}
\end{equation}
where the empirical risk and loss bound are:
\begin{align}
  \Rhat(\lambda) & = \frac{1}{N}\sum_{i=1}^{N}\loss(\lambda; x_i, y_i), &
  b              & = \sup_{\lambda}\,\loss(\lambda; x, y).
\end{align}
This guarantees $\mathbb{E}[\loss(\lamhat; X, Y)] \leq \alpha$ for a new exchangeable point $(X, Y) \sim P$.

\subsection{Label Shift and Importance Weighting}
\label{sec:label_shift}

Under label shift, the deployment distribution $P_{\mathrm{dep}}$ satisfies:
\begin{align}
  \pii{y}{dep}               & \neq \pii{y}{tr},            &
  P_{\mathrm{dep}}(X \mid Y) & = P_{\mathrm{tr}}(X \mid Y).
\end{align}
We parameterise the shift by the odds ratio:
\begin{equation}
  \rhostar = \frac{\pii{1}{dep}\, /\, \pii{0}{dep}}{\pii{1}{tr}\, /\, \pii{0}{tr}}
  \label{eq:rho}
\end{equation}
so $\rhostar > 1$ means the positive class is more prevalent at deployment.
Weighted CRC (W-CRC) restores validity by reweighting calibration losses with importance weights:
\begin{equation}
  w_y = \pii{y}{dep} \,/\, \pii{y}{tr}
  \label{eq:weights}
\end{equation}
as introduced by \citet{tibshirani2019conformal}.
When $\rhostar$ is unknown, Black Box Shift Estimation \citep[BBSE;][]{lipton2018detecting} estimates $w_y$ from the classifier's confusion matrix and unlabelled deployment predictions.

% ======================================================================
\section{Method}
\label{sec:method}

We extend CRC to cost-sensitive decision-making under \emph{uncertain} label shift.
Our method has two components: a tighter excess risk bound (\cref{sec:bound}), and an uncertainty-aware recalibration procedure with three-way decisions that uses this bound to provide guarantees over a shift interval while deferring ambiguous cases (\cref{sec:b4}).

\subsection{Prevalence-Weighted Excess Risk Bound}
\label{sec:bound}

When W-CRC is calibrated with assumed shift $\rhohat$ but the true shift is $\rhostar \neq \rhohat$, the risk guarantee degrades.
The standard approach bounds this degradation using the total variation distance between the assumed and true importance weights, with leading constant \citep{barber2023conformal}:
\begin{equation}
  \Btv = \max(\cfn,\, \cfp).
\end{equation}
We observe that this is loose: the FN and FP contributions to excess risk are each multiplied by the \emph{training prevalence} of the relevant class, which is typically far from~1.

\begin{theorem}[Prevalence-Weighted Bound]
  \label{thm:bound}
  Let $\lamhat$ be calibrated at level $\alpha$ under $P_{\mathrm{tr}}$ with assumed shift $\rhohat$, and let $\rhostar$ be the true shift ratio.
  Define $\boldsymbol{\delta} = \mathbf{w}_{\rhohat} - \mathbf{w}_{\rhostar}$, the importance weight mismatch.
  Then:
  \begin{equation}
    \mathbb{E}_{P_{\mathrm{dep}}}[\loss(\lamhat; X, Y)] \leq \alpha + \Bpw \cdot \lVert \boldsymbol{\delta} \rVert_1 + \frac{\Bpw}{N+1}
    \label{eq:bound}
  \end{equation}
  where $\Bpw = \max\!\bigl(\cfn \cdot \pii{1}{tr},\; \cfp \cdot \pii{0}{tr}\bigr)$.
\end{theorem}

Since $\pii{y}{tr} < 1$ for both classes in any non-degenerate problem, $\Bpw < \Btv$ strictly.
The improvement is largest when classes are imbalanced.
For example, with $\pii{1}{tr} = 0.20$, $\cfn = 1.5$, $\cfp = 1.0$:
\begin{equation}
  \Bpw = \max(0.30,\, 0.80) = 0.80
  \quad\text{vs.}\quad
  \Btv = 1.50
\end{equation}
a 47\% reduction in the leading constant.

\begin{proof}[Proof sketch (full proof in \cref{app:proof})]
  Under label shift, the expected risk decomposes per class:
  \begin{align}
    \Risk_{\rho}(\lambda)
     & = \cfn \cdot \pii{1}{tr} \cdot w_1(\rho) \cdot \mathrm{FNR}(\lambda) \notag        \\
     & \quad + \cfp \cdot \pii{0}{tr} \cdot w_0(\rho) \cdot \mathrm{FPR}(\lambda). \notag
  \end{align}
  The excess $|\Risk_{\rhostar} - \Risk_{\rhohat}|$ is bounded term-by-term, with each term scaled by $c_y \cdot \pii{y}{tr}$ rather than $c_y$ alone.
  The maximum over classes yields $\Bpw \cdot \lVert \boldsymbol{\delta} \rVert_1$; the finite-sample correction adds $\Bpw / (N+1)$.
\end{proof}

The following corollary characterises exactly when the improvement is substantial.

\begin{corollary}[Improvement Regime]
  \label{cor:regime}
  Let $\cfn \geq \cfp$ without loss of generality, and define the crossover prevalence $\pi^* = \cfp / (\cfn + \cfp)$.
  The relative improvement of $\Bpw$ over $\Btv$ is:
  \begin{equation}
    1 - \frac{\Bpw}{\Btv} =
    \begin{cases}
      1 - \pii{1}{tr} & \text{if } \pii{1}{tr} \geq \pi^*, \\[4pt]
      1 - \dfrac{\cfp \cdot \pii{0}{tr}}{\cfn} & \text{if } \pii{1}{tr} < \pi^*.
    \end{cases}
    \label{eq:regime}
  \end{equation}
  In particular, the improvement exceeds $50\%$ when $\pii{1}{tr} < \min\!\bigl(\tfrac{1}{2},\; \cfp / (2\cfn)\bigr)$.
  The tightening is largest when the expensive error corresponds to the rare class---precisely the regime of medical screening, fraud detection, and other cost-sensitive applications with class imbalance.
\end{corollary}

The proof follows directly from the definition of $\Bpw$ and case analysis on which term dominates.

A second consequence of the cost-prevalence interaction explains the counterintuitive empirical finding that uncorrected CRC can become \emph{safer} under positive label shift.

\begin{proposition}[Cost-Asymmetry Monotonicity]
  \label{prop:monotone}
  Under label shift with fixed threshold $\lambda$, the cost-weighted risk
  \begin{equation}
    \Risk(\rho) = \cfn \cdot \pii{1}{dep}(\rho) \cdot \mathrm{FNR}(\lambda) + \cfp \cdot \pii{0}{dep}(\rho) \cdot \mathrm{FPR}(\lambda)
  \end{equation}
  is monotonically decreasing in $\rho$ (for $\rho > 0$) if and only if
  \begin{equation}
    \cfn \cdot \mathrm{FNR}(\lambda) < \cfp \cdot \mathrm{FPR}(\lambda).
    \label{eq:monotone}
  \end{equation}
  When $\cfn > \cfp$, this holds whenever the classifier's operating point satisfies $\mathrm{FNR}(\lambda) / \mathrm{FPR}(\lambda) < \cfp / \cfn$.
\end{proposition}

The proof is immediate: under label shift, FNR and FPR are constant in $\rho$, so $\Risk$ is linear in $\pii{1}{dep}(\rho)$ with slope $\cfn \cdot \mathrm{FNR} - \cfp \cdot \mathrm{FPR}$ (full derivation in \cref{app:prop_proof}).

\Cref{prop:monotone} explains why Standard CRC's risk \emph{decreases} under positive shift in our LLM toxicity experiment (\cref{sec:results}): the calibrated threshold satisfies $\mathrm{FNR}/\mathrm{FPR} < \cfp/\cfn$, so increasing $\rho$ actually reduces risk.
This ``accidentally safe'' behaviour reverses under negative shift ($\rho < 1$), where Standard CRC violates in $97\%$ of seeds.

\subsection{Uncertainty-Aware Recalibration and Three-Way Decisions}
\label{sec:b4}

Given an uncertainty interval $[\rho_{\mathrm{lo}}, \rho_{\mathrm{hi}}]$ on the shift ratio, we recalibrate CRC to guarantee risk control for \emph{all} $\rhostar$ in that interval, and let deferral emerge from endpoint disagreement.

\textbf{Recalibration.}
\label{sec:bayesian}
Using \cref{thm:bound}, the worst-case excess risk over the interval is:
\begin{align}
  \epsilon_{\max}
   & = \Bpw \cdot \max\!\bigl(
  \lVert \boldsymbol{\delta}_{\rho_{\mathrm{lo}}} \rVert_1,\,
  \lVert \boldsymbol{\delta}_{\rho_{\mathrm{hi}}} \rVert_1
  \bigr) + \frac{\Bpw}{N+1}
  \label{eq:excess}
\end{align}
since the excess is monotone in $|\rho - \rhohat|$.
We calibrate CRC at the adjusted level $\asafe = \alpha - \epsilon_{\max}$ instead of $\alpha$, guaranteeing $\mathbb{E}_{P_{\mathrm{dep}}}[\loss(\lamhat; X, Y)] \leq \alpha$ for any $\rhostar \in [\rho_{\mathrm{lo}}, \rho_{\mathrm{hi}}]$.
If $\asafe \leq 0$, the interval is too wide for safe calibration; the practitioner must narrow it.
The key tradeoff is: wider interval $\Rightarrow$ $\epsilon_{\max}\!\uparrow$ $\Rightarrow$ $\asafe\!\downarrow$ $\Rightarrow$ higher FNR.

\textbf{Obtaining the interval.}
The interval $[\rho_{\mathrm{lo}}, \rho_{\mathrm{hi}}]$ can be specified in two ways:
\begin{itemize}
  \item \emph{Label-free mode:} the practitioner specifies it from domain knowledge (e.g., $[0.8, 1.2]$). No deployment data is needed.
  \item \emph{Data-driven mode:} when unlabelled deployment data is available, we estimate the interval via a Beta-Binomial model. BBSE provides a point estimate $\hat{\pi}_{\mathrm{dep}}$; we convert this to pseudo-observations and compute the posterior:
  \begin{equation}
    \pi_{\mathrm{dep}} \mid \text{data} \;\sim\; \mathrm{Beta}(a + n_+,\; b + M - n_+)
    \label{eq:posterior}
  \end{equation}
  with default prior $a = b = 2$.  The $95\%$ credible interval on $\rho$ (obtained by Monte Carlo transform) replaces the hand-specified interval.
\end{itemize}

\textbf{Three-way decisions.}
\label{sec:threeway}
Rather than forcing a binary decision, we calibrate CRC at \emph{both} interval endpoints, producing thresholds $\lambda^{\mathrm{lo}}$ and $\lambda^{\mathrm{hi}}$.
The decision rule is:
\begin{equation}
  f(x) < \lambda_{\min} \Rightarrow \text{negative};\quad
  f(x) \geq \lambda_{\max} \Rightarrow \text{positive};\quad
  \text{otherwise} \Rightarrow \textbf{defer} \notag
\end{equation}
where $\lambda_{\min} = \min(\lambda^{\mathrm{lo}}, \lambda^{\mathrm{hi}})$ and $\lambda_{\max} = \max(\lambda^{\mathrm{lo}}, \lambda^{\mathrm{hi}})$.
The defer zone $[\lambda_{\min}, \lambda_{\max})$ contains scores where the two endpoint calibrations disagree.
More shift uncertainty $\Rightarrow$ wider defer zone $\Rightarrow$ more deferral.
When $\rho_{\mathrm{lo}} = \rho_{\mathrm{hi}}$, $\lambda_{\min} = \lambda_{\max}$ and the method reduces to standard binary CRC.
No free parameters are introduced.

\begin{theorem}[Conditional Risk Guarantee]
  \label{thm:threeway}
  For any $\rhostar \in [\rho_{\mathrm{lo}}, \rho_{\mathrm{hi}}]$, let $d(\rhostar) = P_{\rhostar}(\lambda_{\min} \leq f(X) < \lambda_{\max})$ be the deferral rate.  Then:
  \begin{equation}
    \mathbb{E}_{\rhostar}\bigl[\loss_{\mathrm{3way}}(X, Y) \;\big|\; X \in \mathrm{decided}\bigr]
    \;\leq\; \frac{\alpha}{1 - d(\rhostar)}.
    \label{eq:threeway}
  \end{equation}
  Moreover, $d(\rhostar)$ is observable from unlabelled deployment scores.
\end{theorem}

\begin{proof}
  For $j \in \{\mathrm{lo}, \mathrm{hi}\}$, W-CRC gives $\mathbb{E}_{\rho_j}[\loss(\lambda^j; X, Y)] \leq \alpha$.
  For any $\rhostar$ in the interval, the excess at $\rhostar$ does not exceed the excess at the farther endpoint (by monotonicity of $\lVert \boldsymbol{\delta} \rVert_1$ in $|\rho - \rhohat|$), so $\mathbb{E}_{\rhostar}[\loss(\lambda^j)] \leq \alpha$.
  On the decided set $D$, the three-way loss equals the single-threshold loss under $\lambda^j$ (scores above $\lambda_{\max}$ or below $\lambda_{\min}$ are classified identically by both thresholds).
  Decomposing over decided and deferred sets: $(1 - d) \cdot \mathbb{E}_{\rhostar}[\loss_{\mathrm{3way}} \mid D] \leq \alpha$.
\end{proof}

% ======================================================================
\section{Experimental Setup}
\label{sec:setup}

\subsection{Tasks}
\label{sec:tasks}

\textbf{LLM Toxicity Detection (headline).}
Binary toxicity classification on the CivilComments dataset \citep{borkan2019nuanced}, scored by OpenAI's moderation API (harassment category score).
Calibration pool: 10{,}000 general comments with low identity-attack score (${\sim}6\%$ toxic).
Deployment pool: 10{,}000 identity-targeting comments with identity-attack score $\geq 0.1$ (${\sim}29\%$ toxic), creating natural label shift ($\rhostar \approx 6.3$).
Cost asymmetry $\cfn = 5$, $\cfp = 1$: missing toxic content directed at identity groups is far costlier than over-blocking.
Toxicity subcategory groups (Insult, Threat, Obscene, Identity Attack, Severe) enable per-group burden analysis.
Label shift is simulated via rejection sampling at $\rhostar \in \{0.5, 1.0, 1.5, 2.0, 3.0, 5.0\}$.
This experiment directly addresses the central LLM safety question: \emph{when should the system escalate to a human moderator?}

\textbf{Folktables ACSIncome.}
Income $>\$50$k prediction on the 2018 American Community Survey \citep{fisch2022conformal}.
  A HistGradientBoosting classifier (AUROC 0.90) is trained on five low-prevalence states (MS, AR, AL, WV, LA) and deployed on five high-prevalence states (MD, NJ, CT, MA, VA).
  A three-way split---70\% model training, 30\% calibration pool, separate deploy states---ensures calibration scores are out-of-sample.
  Label shift is simulated via rejection sampling at $\rhostar \in \{0.2, 0.5, 0.67, 1.0, 1.5, 2.0, 3.0, 5.0\}$, spanning both negative shift ($\rhostar < 1$, disease becomes rarer) and positive shift ($\rhostar > 1$, disease becomes more prevalent).

  \textbf{GAIA Agentic QA.}
  Open-ended question answering on the GAIA benchmark \citep{mialon2023gaia}, a 165-task dataset spanning three difficulty levels (L1: 53, L2: 86, L3: 26) with 23\% requiring file attachments.
  A GPT-4o-mini agent with simulated tool access (web search, calculator) generates first-pass answers; an LLM-as-judge produces \emph{badness scores} $s \in [0, 1]$ estimating $P(\text{answer is wrong})$.
  Calibration: L1 tasks (easier).
  Deployment: L2+L3 tasks (harder), creating \emph{failure-prevalence shift}---the agentic analogue of label shift.
  The three-way decision rule maps badness scores to actions: answer ($s < \lambda_{\text{lo}}$), retry with reflection ($\lambda_{\text{lo}} \leq s < \lambda_{\text{hi}}$), or hand off to a human ($s \geq \lambda_{\text{hi}}$).
  This experiment extends the CRC framework from classifier outputs to agentic decision gating.

  \subsection{Baselines}
  \label{sec:baselines}

  We compare five methods, all using the cost-weighted loss (\cref{eq:loss}):
  \begin{itemize}
    \item \textbf{Standard CRC}: no shift correction; calibrates at $\alpha$.
    \item \textbf{Oracle W-CRC}: importance-weighted CRC with true $\rhostar$ (upper bound on what shift correction can achieve).
    \item \textbf{BBSE W-CRC}: importance-weighted CRC with BBSE-estimated weights.
    \item \textbf{Ours (label-free)}: our method in label-free mode; $\rhohat = 1.0$, $[\rho_{\mathrm{lo}}, \rho_{\mathrm{hi}}] = [0.8, 1.2]$.
    \item \textbf{Ours (data-driven)}: our method in data-driven mode with three-way decisions; $\mathrm{Beta}(2, 2)$ prior, $95\%$ credible level.
  \end{itemize}
  For the GAIA agentic experiment, we use the same four gating policies (Always Answer, Standard CRC, Ours label-free, Ours data-driven) but map decisions to agent actions: \emph{answer} (respond autonomously), \emph{retry} (re-attempt with reflection prompt), or \emph{handoff} (escalate to human).

  \subsection{Metrics}
  \label{sec:metrics}

  For each method and shift level we report the following on $M$ deployment samples over 30 random seeds (mean $\pm$ std):
  \begin{itemize}
    \item \textbf{Cost-weighted risk:}
          $\displaystyle\Rhat = \frac{1}{M}\sum_{j=1}^{M}\loss(\lamhat; x_j, y_j)$
    \item \textbf{Violation rate:} fraction of seeds where $\Rhat > \alpha$
    \item \textbf{FNR / FPR:} to disaggregate cost components
    \item \textbf{Deferral rate:} fraction of samples deferred (Ours data-driven only)
    \item \textbf{Risk on decided:} cost-weighted risk computed only on non-deferred samples
  \end{itemize}

  \subsection{Computational Details}
  \label{sec:details}

  Shared parameters across all experiments:
  \begin{align}
    \alpha    & = 0.30,          &
    N         & = 1{,}000,       &
    M         & = 2{,}000,       &
    |\Lambda| & = 1{,}000 \notag
  \end{align}
  where $|\Lambda|$ is the threshold grid size over $[0, 1]$.
  Cost parameters: $\cfn = 5.0$, $\cfp = 1.0$ for LLM toxicity (missing toxic content is 5$\times$ costlier than over-blocking); $\cfn = 1.5$, $\cfp = 1.0$ for Folktables.
  Label-free recalibration interval: $[0.8, 1.2]$ (approximate exchangeability with 20\% uncertainty).
  LLM toxicity: 20{,}000 CivilComments scored by OpenAI's \texttt{omni-moderation-latest} model (harassment category score).
  Folktables: HistGradientBoosting (\texttt{max\_iter}${}=200$, \texttt{max\_depth}${}=6$, \texttt{lr}${}=0.1$).

  % ======================================================================
  \section{Results}
  \label{sec:results}

  We organise results around safety-critical LLM deployment, synthesising evidence across datasets.
  The two main LLM-facing settings are toxicity moderation and GAIA agentic QA; Folktables provides non-LLM support for the same risk-control mechanisms.

  % --- Take-away 0: LLM headline ---
  \textbf{On LLM toxicity detection, cost asymmetry amplifies risk under negative shift, and our label-free mode provides the widest safety margin.}
  \Cref{tab:llm,fig:llm_risk} show cost-weighted risk on CivilComments ($\cfn = 5$, $\cfp = 1$, 30 seeds).
  With high cost asymmetry, Standard CRC violates at $\rhostar = 0.5$ (97\% of seeds) and $\rhostar = 1.0$ (37\%)---the counterintuitive direction predicted by \cref{prop:monotone}.
  Ours (label-free) eliminates all violations across every shift level by calibrating at $\asafe = 0.228$ (risk $0.225$ at $\rhostar = 1.0$ vs.\ $0.292$ for Standard CRC---a $23\%$ safety margin).
  Ours (data-driven) matches Oracle and BBSE correction quality while deferring ${\sim}1.4\%$ of ambiguous cases to human review, answering ``when to escalate'' with no tuning.
  \Cref{fig:llm_subgroup} reveals that the false-negative burden is distributed unevenly across toxicity subcategories, with smaller groups bearing disproportionate FNR.

  \begin{table}[t]
    \caption{Cost-weighted risk on CivilComments LLM toxicity under label shift ($\alpha = 0.30$, $\cfn = 5$, $\cfp = 1$, 30 seeds). Stars ($*$) denote violation ($\Rhat > \alpha$ in $> 0\%$ of seeds). \textbf{Take-away:} Ours (label-free) eliminates violations everywhere; Standard CRC fails under negative shift.}
    \label{tab:llm}
    \begin{center}
      \begin{small}
        \input{tables/llm_table.tex}
      \end{small}
    \end{center}
    \vskip -0.1in
  \end{table}

  \begin{figure}[t]
    \centering
    \includegraphics[width=\columnwidth]{figures/fig_llm_risk_vs_rho.png}
    \caption{Cost-weighted risk on LLM toxicity detection as $\rhostar$ increases. \textbf{Top:} Mean risk with $\pm 1\sigma$ bands; the high cost asymmetry ($\cfn/\cfp = 5$) amplifies violations under shift. \textbf{Bottom:} Per-seed violation rate.}
    \label{fig:llm_risk}
  \end{figure}

  \begin{figure}[t]
    \centering
    \includegraphics[width=\columnwidth]{figures/fig_llm_subgroup_burden.png}
    \caption{Per-identity-group FNR on LLM toxicity at $\rhostar = 2.0$. \textbf{Take-away:} aggregate risk control masks unequal burden; identity groups with higher baseline toxicity prevalence bear disproportionate false-negative rates.}
    \label{fig:llm_subgroup}
  \end{figure}

  % --- Take-away 1 ---
  \textbf{Our label-free mode provides the widest safety margin without oracle access, at the cost of increased conservatism.}
  \Cref{tab:main} shows cost-weighted risk across shift levels on Folktables (30 seeds).
  Under exchangeability ($\rhostar = 1.0$), all methods control risk, but ours (label-free) provides the largest margin: risk $0.221$ vs.\ $0.259$ for Standard CRC, by calibrating at $\asafe = 0.246$ rather than $\alpha = 0.30$.
  This margin persists under moderate shift ($\rhostar = 1.5$): ours (label-free) risk $0.236$ vs.\ $0.249$ for Standard CRC.

  \begin{table}[t]
    \caption{Cost-weighted risk on Folktables under simulated label shift ($\alpha = 0.30$, 30 seeds). \textbf{Take-away:} all methods control risk at $\rhostar = 1.0$, but ours (label-free) provides the widest safety margin (risk $0.221$ vs.\ $0.259$); it degrades beyond its assumed interval at $\rhostar \geq 2.0$.}
    \label{tab:main}
    \begin{center}
      \begin{small}
        \begin{tabular}{llcc}
          \toprule
          $\rhostar$ & Method       & Risk                      & Viol.\       \\
          \midrule
          1.0        & Standard CRC & $0.259 \pm 0.02$          & 0\%          \\
          1.0        & BBSE W-CRC   & $0.254 \pm 0.02$          & 0\%          \\
          1.0        & Ours (l-free) & $\mathbf{0.221 \pm 0.01}$ & \textbf{0\%} \\
          \midrule
          2.0        & Standard CRC & $0.237 \pm 0.01$          & 0\%          \\
          2.0        & Oracle W-CRC & $0.270 \pm 0.02$          & 3\%          \\
          2.0        & BBSE W-CRC   & $0.259 \pm 0.02$          & 0\%          \\
          2.0        & Ours (l-free) & $0.254 \pm 0.02$          & 7\%          \\
          \midrule
          5.0        & Standard CRC & $0.209 \pm 0.02$          & 0\%          \\
          5.0        & BBSE W-CRC   & $0.264 \pm 0.03$          & 13\%         \\
          5.0        & Ours (l-free) & $0.303 \pm 0.05$          & 53\%         \\
          \bottomrule
        \end{tabular}
      \end{small}
    \end{center}
    \vskip -0.1in
  \end{table}

  \begin{figure}[t]
    \centering
    \includegraphics[width=\columnwidth]{figures/fig_real_risk_vs_rho.png}
    \caption{Cost-weighted risk on Folktables as $\rhostar$ increases. \textbf{Top:} Mean risk with $\pm 1\sigma$ bands across seeds; Standard CRC's upper band crosses $\alpha$ at moderate shift despite safe-looking means. \textbf{Bottom:} Per-seed violation rate; Standard CRC violates $\alpha$ in ${\sim}40\%$ of seeds at $\rho = 0.2$, confirming that mean risk masks individual failures.}
    \label{fig:risk_vs_rho}
  \end{figure}

  % --- Take-away 2 ---
  \textbf{Cost asymmetry produces counterintuitive risk dynamics: uncorrected CRC becomes accidentally safe under positive shift but \emph{violates} under negative shift (\cref{prop:monotone}).}
  \Cref{fig:risk_vs_rho} and \cref{tab:main} show this clearly: Standard CRC's risk \emph{decreases} from $0.259$ at $\rhostar = 1.0$ to $0.209$ at $\rhostar = 5.0$ (positive shift), but \emph{increases} to $0.295$ at $\rhostar = 0.2$ (negative shift), violating the $\alpha = 0.30$ guarantee in $40\%$ of seeds.
  \Cref{prop:monotone} identifies the exact condition: $\cfn \cdot \mathrm{FNR}(\lambda) < \cfp \cdot \mathrm{FPR}(\lambda)$, which holds at our calibrated threshold ($\mathrm{FNR}/\mathrm{FPR} \approx 0.34 < \cfp/\cfn = 0.67$).
  The accidentally-safe effect operates only under positive shift; negative shift reverses the inequality and exposes the unprotected direction.

  % --- Take-away 3 ---
  \textbf{The tighter bound translates to concrete utility gains and extends recalibration feasibility.}
  \Cref{tab:utility} compares our label-free recalibration with $\Bpw$ vs.\ $\Btv$ across uncertainty intervals of increasing width.
  At $[0.6, 1.5]$, the tighter bound yields $43\%$ more positive predictions ($37.8\%$ vs.\ $26.4\%$).
  At $[0.5, 2.0]$, the TV bound makes recalibration \emph{infeasible} ($\asafe < 0$), while $\Bpw$ still permits safe calibration ($\asafe = 0.086$).
  The maximum feasible interval width is $63\%$ larger with $\Bpw$.
  Synthetic validation (\cref{app:additional}) confirms the bound is 27--39\% tighter across $\rhostar \in \{2, 5, 10\}$; as \cref{cor:regime} predicts, the gain is largest when the expensive error ($\cfn$) corresponds to the rare class.

  \begin{table}[t]
    \caption{Decision-quality comparison: label-free recalibration with $\Bpw$ vs.\ $\Btv$ bound across uncertainty intervals ($\alpha = 0.30$, $N = 1{,}000$, 20 seeds). \textbf{Take-away:} the tighter bound yields higher $\asafe$, more positive predictions, and extends recalibration feasibility to wider intervals.}
    \label{tab:utility}
    \begin{center}
      \begin{small}
        \begin{tabular}{lcccccc}
          \toprule
          Interval & \multicolumn{2}{c}{$\asafe$} & \multicolumn{2}{c}{Pos.\ rate} & \multicolumn{2}{c}{FNR} \\
          \cmidrule(lr){2-3} \cmidrule(lr){4-5} \cmidrule(lr){6-7}
                   & $\Bpw$ & $\Btv$ & $\Bpw$ & $\Btv$ & $\Bpw$ & $\Btv$ \\
          \midrule
          $[0.8, 1.2]$ & $.246$ & $.199$ & $44.0\%$ & $39.3\%$ & $0.2\%$ & $0.4\%$ \\
          $[0.6, 1.5]$ & $.183$ & $.080$ & $37.8\%$ & $26.4\%$ & $0.5\%$ & $2.7\%$ \\
          $[0.5, 2.0]$ & $.086$ & --- & $27.0\%$ & --- & $2.4\%$ & --- \\
          \bottomrule
        \end{tabular}
      \end{small}
    \end{center}
    \vskip -0.1in
  \end{table}

  % --- Take-away 4 ---
  \textbf{Three-way decisions provide guaranteed risk control on the decided subset (\cref{thm:threeway}), with zero free parameters.}
  \Cref{fig:deferral} shows the deferral curve from Experiment~7: as the Bayesian prior becomes less informative ($k \to 0$, i.e., wider credible interval), the deferral rate increases smoothly from near $0\%$ (strong prior, $k = 50$) to the maximum deferral level (weak prior, $k = 0.5$).
  On Folktables at $\rhostar = 2.0$, ours (data-driven) maintains risk control on the decided subset (\cref{thm:threeway}: conditional risk $\leq \alpha / (1 - d)$) while deferring the ambiguous cases---those where the shift interval creates disagreement about the safe threshold.
  Critically, no deferral cost hyperparameter is needed: the defer zone is entirely determined by the width of the Bayesian credible interval.

  \begin{figure}[t]
    \centering
    \includegraphics[width=\columnwidth]{figures/fig7_deferral_curve.png}
    \caption{Deferral rate vs Bayesian prior strength (synthetic, $\rhostar = 3.0$). \textbf{Take-away:} uninformative priors (small $k$) produce wide shift intervals and high deferral; informative priors (large $k$) narrow the interval and reduce deferral. The defer zone emerges from shift uncertainty with no free parameters.}
    \label{fig:deferral}
  \end{figure}

  % --- GAIA Agentic ---
  \subsection{Agentic Decision Gating on GAIA}
  \label{sec:gaia_results}

  \textbf{CRC-based gating reduces autonomous risk by 47\% under task-difficulty shift.}
  We apply the CRC framework to LLM agent decisions on the GAIA benchmark, where the ``score'' is a badness estimate $s \in [0,1]$ from an LLM judge and the ``label'' is whether the agent's answer is wrong.
  \Cref{tab:gaia} and \cref{fig:gaia} show results under difficulty shift (calibrate on L1, deploy on L2+L3, 40 deployment tasks).
  Without gating (Always Answer), the agent produces wrong answers on $37.5\%$ of deployment tasks.
  All CRC-based policies reduce autonomous risk to $20.0\%$---a $47\%$ reduction---by routing $62.5\%$ of tasks to human operators.
  The $30\%$ of tasks answered autonomously achieve a no-human success rate of $30\%$, with the remaining risk concentrated in the hardest tasks where the agent's confidence is highest but incorrect.

  Under modality shift (calibrate on text-only, deploy on attachment-requiring tasks), CRC gating eliminates autonomous risk entirely: all attachment tasks are routed to handoff, correctly identifying that the agent lacks the capability to process file attachments.

  \textbf{When the agent is out of its depth, the method safely escalates 100\% of tasks.}
  On the full 165-task real GAIA benchmark (where the agent achieves only $2.4\%$ accuracy due to simulated rather than live tools), all CRC-based policies hand off $100\%$ of deployment tasks.
  This is the correct behaviour: the badness scores are uniformly high, the calibrated threshold reflects near-total failure, and the method refuses to answer autonomously.
  This ``degenerate but correct'' outcome demonstrates that CRC gating fails safely---it does not produce false confidence when the underlying agent is incapable.

  \begin{table}[t]
    \caption{GAIA agentic decision gating under task-difficulty shift (L1 $\to$ L2+L3, 40 deployment tasks). \textbf{Take-away:} CRC-based gating reduces autonomous risk from $37.5\%$ to $20.0\%$ by routing uncertain tasks to human operators.}
    \label{tab:gaia}
    \begin{center}
      \begin{small}
        \input{tables/gaia_table.tex}
      \end{small}
    \end{center}
    \vskip -0.1in
  \end{table}

  \begin{figure}[t]
    \centering
    \includegraphics[width=\columnwidth]{figures/fig_gaia_policy.png}
    \caption{GAIA policy comparison under difficulty shift.  \textbf{Left:} autonomous risk (fraction of wrong autonomous answers). \textbf{Centre:} autonomous answer rate. \textbf{Right:} handoff + retry rate.  CRC-based policies trade answer rate for safety, routing $62.5\%$ of tasks to human review.}
    \label{fig:gaia}
  \end{figure}

  % ======================================================================
  \section{Discussion}
  \label{sec:discussion}

  \textbf{When the method works and when it fails.}
  Our recalibration provides guaranteed cost-weighted risk control under two conditions: (i)~the true shift $\rhostar$ falls within the interval, and (ii)~the label-shift assumption $P(X \mid Y) = \text{const}$ holds.
  When violated---e.g., $\rhostar$ outside the interval (Folktables, $\rhostar = 5.0$)---the guarantee breaks, but the failure mode is transparent and testable.
  The three-way extension answers the operational question \emph{when should a deployed LLM say ``I don't know''?}---on LLM toxicity, the conditional risk on decided cases is at most $\alpha / (1 - d)$ (\cref{thm:threeway}), with $d$ observable from unlabelled scores.
  The deferral rate doubles as a staffing metric: it tells the practitioner how many cases require human review.

  \textbf{From classifiers to agents: failure-prevalence as label shift.}
  The GAIA experiment demonstrates that CRC extends naturally from classifier confidence scores to agentic badness scores.
  The key insight is that \emph{failure prevalence} plays the role of class prevalence: when deployment tasks are harder than calibration tasks, the fraction of wrong answers shifts upward, creating a failure-prevalence shift analogous to label shift.
  The method works well under controlled conditions (demo tasks with $68\%$ first-pass accuracy) but becomes maximally conservative when the agent fundamentally lacks capability ($2.4\%$ accuracy on real GAIA).
  In LLM-safety terms, this is precisely the desired failure mode: abstain rather than hallucinate with confidence.

  \textbf{Limitations.}
  (1)~The LLM toxicity scorer (OpenAI's moderation API) is a black-box; score calibration quality depends on the API model version.
  (2)~The Bayesian interval relies on BBSE's point estimate; while coverage remains $\geq 84\%$ across all tested shift levels, ill-conditioned confusion matrices could further degrade it.
  (3)~Our bound, while tighter than the TV bound and yielding concrete utility gains (\cref{tab:utility}), remains conservative relative to realised risk.
  (4)~We evaluate binary classification only; extension to multi-class toxicity categories with structured costs is non-trivial.
  (5)~The GAIA agent uses simulated tools (no live web search), limiting first-pass accuracy to $2.4\%$ on real tasks; with live tools the method would operate in a more informative regime where the answer/retry/handoff distinctions become meaningful at scale.

  \textbf{Future work.}
  While \cref{thm:threeway} provides a conditional risk guarantee of $\alpha / (1-d)$, tighter bounds that exploit the score distribution within the decided region may be achievable.
  For LLM safety, extending the framework to multi-label toxicity detection (harassment, hate speech, self-harm) with structured cost matrices is a natural next step.
  Beyond this: online conformal risk control for streaming content moderation, using LLM logprobs as native confidence scores for hallucination prevention and agentic deferral, and scaling the gating layer to multi-step agent workflows with live tool access where intermediate failure modes (tool errors, reasoning loops) create richer score distributions.

  % ======================================================================
  \section{Conclusion}
  \label{sec:conclusion}

  We have extended Conformal Risk Control to cost-sensitive decision-making under uncertain label shift, motivated by a practical LLM-safety question: \emph{when should a model answer, and when should it say ``I don't know''?}
  Our prevalence-weighted bound tightens the standard TV bound by 27--39\%, and the recalibration procedure translates uncertainty about shift magnitude into an explicit safety--utility tradeoff---with three-way abstention emerging from endpoint disagreement with no free parameters (\cref{thm:threeway}).
  On LLM toxicity detection, recalibration eliminates all risk violations while three-way decisions answer ``when to escalate'' with formal guarantees.
  On the GAIA agentic benchmark, CRC-based gating reduces autonomous risk in a controlled difficulty-shift demo and correctly abstains when the underlying agent lacks capability on the full benchmark.
  The main message is deployment-time rather than training-time: we do not remove hallucinations from the model, but we can make a deployed LLM far less likely to answer unsafely when uncertainty is high.

  % ======================================================================
  \section*{Impact Statement}

  This paper studies a deployment-time safety layer for LLM systems under distribution shift, especially in settings where a confident wrong answer is worse than abstaining.
  Our primary finding is methodological: risk control can be used to decide when a deployed LLM should answer, retry, or hand off, but natural deployment shift can still violate the underlying assumptions and subgroup burden can remain hidden beneath aggregate metrics.
  The LLM toxicity experiment shows that subgroup burden can remain uneven even when aggregate risk is controlled; the GAIA experiment shows that the same machinery can force abstention when an agent is out of its depth.
  These findings are demonstrated on public benchmarks, not in production deployments.
  Any application to production medical, legal, or other safety-critical decision support would require prospective validation with held-out scoring models, domain experts, and human-adjudicated outcomes, which is beyond the scope of this work.
  Our results should therefore be read as a proof of concept for safer LLM deployment, not as evidence of deployment readiness in high-stakes domains.

  \bibliography{references}
  \bibliographystyle{icml2026}

  %%%%%%%%%%%%%%%%%%%%%%%%%%%%%%%%%%%%%%%%%%%%%%%%%%%%%%%%%%%%%%%%%%%%%%%%%%%%%%%
  \newpage
  \appendix
  \onecolumn

  \section{Proof of \texorpdfstring{\cref{thm:bound}}{Theorem 4.1}}
  \label{app:proof}

  Let $\Risk_{\rho}(\lambda) = \mathbb{E}_{P_{\rho}}[\loss(\lambda; X, Y)]$ denote the expected cost-weighted risk under shift ratio $\rho$.
  Under label shift, $P_{\mathrm{dep}}(X \mid Y) = P_{\mathrm{tr}}(X \mid Y)$, so the risk decomposes as:
  \begin{align}
    \Risk_{\rho}(\lambda)
     & = \cfn \cdot \pii{1}{tr} \cdot w_1(\rho) \cdot \mathrm{FNR}(\lambda)
    + \cfp \cdot \pii{0}{tr} \cdot w_0(\rho) \cdot \mathrm{FPR}(\lambda)
    \label{eq:risk_decomp}
  \end{align}
  where $w_y(\rho) = P_\rho(Y = y) / \pii{y}{tr}$ and $\mathrm{FNR}(\lambda) = P(f(X) < \lambda \mid Y = 1)$, $\mathrm{FPR}(\lambda) = P(f(X) \geq \lambda \mid Y = 0)$ are properties of the scoring function under $P_{\mathrm{tr}}(X \mid Y)$.

  The difference between the risk under the true shift $\rhostar$ and the assumed shift $\rhohat$ is:
  \begin{align}
    |\Risk_{\rhostar}(\lambda) - \Risk_{\rhohat}(\lambda)|
     & = \bigl|\cfn \cdot \pii{1}{tr} \cdot \bigl(w_1(\rhostar) - w_1(\rhohat)\bigr) \cdot \mathrm{FNR}(\lambda)
    + \cfp \cdot \pii{0}{tr} \cdot \bigl(w_0(\rhostar) - w_0(\rhohat)\bigr) \cdot \mathrm{FPR}(\lambda)\bigr|    \\
     & \leq \cfn \cdot \pii{1}{tr} \cdot |w_1(\rhostar) - w_1(\rhohat)|
    + \cfp \cdot \pii{0}{tr} \cdot |w_0(\rhostar) - w_0(\rhohat)|                                                \\
     & \leq \Bpw \cdot \bigl(|w_1(\rhostar) - w_1(\rhohat)| + |w_0(\rhostar) - w_0(\rhohat)|\bigr)               \\
     & = \Bpw \cdot \lVert \boldsymbol{\delta} \rVert_1
  \end{align}
  where $\Bpw = \max(\cfn \cdot \pii{1}{tr},\, \cfp \cdot \pii{0}{tr})$ and $\boldsymbol{\delta} = \mathbf{w}_{\rhohat} - \mathbf{w}_{\rhostar}$.

  CRC calibrated at level $\alpha$ under $\rhohat$ guarantees:
  \begin{equation}
    \Risk_{\rhohat}(\lamhat) \leq \alpha + \frac{\Bpw}{N+1}
  \end{equation}
  Combining:
  \begin{equation}
    \Risk_{\rhostar}(\lamhat) \leq \Risk_{\rhohat}(\lamhat) + \Bpw \cdot \lVert \boldsymbol{\delta} \rVert_1 \leq \alpha + \Bpw \cdot \lVert \boldsymbol{\delta} \rVert_1 + \frac{\Bpw}{N+1}
  \end{equation}
  which is \cref{eq:bound}. \qed

  \section{Proof of \texorpdfstring{\cref{prop:monotone}}{Proposition 4.3}}
  \label{app:prop_proof}

  Under label shift, $P(X \mid Y = y)$ is invariant, so $\mathrm{FNR}(\lambda)$ and $\mathrm{FPR}(\lambda)$ are constant in $\rho$.
  The deployment prevalence $\pii{1}{dep}(\rho) = \rho \cdot \pii{1}{tr} / (\rho \cdot \pii{1}{tr} + \pii{0}{tr})$ is strictly increasing in $\rho$, and $\pii{0}{dep} = 1 - \pii{1}{dep}$.
  Since $\Risk$ is linear in $\pii{1}{dep}$, $\partial \Risk / \partial \pii{1}{dep} = \cfn \cdot \mathrm{FNR}(\lambda) - \cfp \cdot \mathrm{FPR}(\lambda)$.
  This is negative precisely when $\cfn \cdot \mathrm{FNR}(\lambda) < \cfp \cdot \mathrm{FPR}(\lambda)$. \qed

  \section{Experimental Details}
  \label{app:details}

  \subsection{LLM Toxicity Detection (CivilComments)}

  Data from the CivilComments dataset \citep{borkan2019nuanced} via HuggingFace \texttt{datasets}.
  20{,}000 comments (10{,}000 calibration pool, 10{,}000 deployment pool).
  Calibration pool: comments with identity-attack score $< 0.1$ ($6.0\%$ toxic).
  Deployment pool: comments with identity-attack score $\geq 0.1$ ($28.7\%$ toxic), creating natural label shift ($\rhostar \approx 6.3$).
  Scoring: OpenAI's \texttt{omni-moderation-latest} moderation API, \texttt{harassment} category score $\in [0, 1]$.
  Ground truth: binarised at toxicity $\geq 0.5$.
  Cost parameters: $\cfn = 5.0$ (missing toxic content), $\cfp = 1.0$ (over-blocking).
  Subgroups by dominant toxicity subcategory: Insult (56.2\%), Identity Attack (37.5\%), Threat (4.1\%), Obscene (2.1\%), Severe ($<0.1\%$).
  Label shift simulated via rejection sampling at $\rhostar \in \{0.5, 1.0, 1.5, 2.0, 3.0, 5.0\}$.

  \subsection{Folktables ACSIncome}

  Data from the 2018 American Community Survey via the \texttt{folktables} package.
  Training states (low positive prevalence): MS, AR, AL, WV, LA.
  Deployment states (high positive prevalence): MD, NJ, CT, MA, VA.
  Classifier: \texttt{HistGradientBoostingClassifier} with \texttt{max\_iter=200}, \texttt{max\_depth=6}, \texttt{learning\_rate=0.1}, \texttt{random\_state=42}.
  Three-way split: model training (70\%), calibration pool (30\%), separate deployment states.
  AUROC: 0.90 (training), 0.88 (calibration pool).
  Label shift simulated via rejection sampling: positive samples retained with probability $\min(1, \rhostar \cdot \pii{1}{tr} / \pii{0}{tr})$, then randomly downsampled to $N$ calibration and $M$ deployment points.

  \subsection{Folktables ACSPublicCoverage}

  Data from the 2018 ACS using a disjoint set of states from ACSIncome.
  Training states (Medicaid non-expansion, low coverage prevalence 21.5\%): TX, FL, GA, NC, TN.
  Deployment states (Medicaid expansion, high coverage prevalence 36.9\%): NY, CA, MN, OR, WA.
  Classifier: same HistGradientBoosting configuration as ACSIncome.
  Three-way split: identical protocol (70\% model training, 30\% calibration pool, separate deployment states).
  Natural shift ratio: $\rhostar \approx 2.1$, driven by Medicaid policy rather than geographic income variation.

  \subsection{GAIA Agentic QA}

  Data from the GAIA benchmark \citep{mialon2023gaia} via HuggingFace (gated dataset, validation split).
  165 tasks: 53 Level~1, 86 Level~2, 26 Level~3; 38 (23\%) require file attachments.
  Agent: GPT-4o-mini with tool-calling (simulated web search and calculator; no live tool access).
  First-pass accuracy: $2.4\%$ on real tasks (limited by simulated tools); $68\%$ on controlled demo tasks.
  Scoring: LLM-as-judge (GPT-4o-mini) produces badness scores $s \in [0,1]$.
  Difficulty split: L1 calibration, L2+L3 deployment (40 tasks).
  Modality split: text-only calibration (127 tasks), attachment-requiring deployment (38 tasks).
  Demo experiment uses 60 controlled tasks (20 per level) with natural difficulty gradient and known ground truth.
  Cost parameters: $\cfn = 5.0$ (wrong autonomous answer), $\cfp = 1.0$ (unnecessary handoff).

  \subsection{PACPS Heart Disease Replication}

  CDC Heart Disease dataset, following the protocol of \citet{fisch2022conformal}.
  Methods: PS (uncorrected), PS-C (oracle correction), PS-R (point estimate), PS-W (worst-case), PS-MIW (max importance weight).
  3 seeds.
  PS violates coverage in 100\% of runs; all corrected methods achieve 0\% violation rate, matching the published results.

  \section{Additional Results}
  \label{app:additional}

  \subsection{Bound Tightness}

  \begin{table}[h]
    \caption{Prevalence-weighted bound vs.\ TV bound across shift ratios ($\pii{1}{tr} = 0.20$, $\cfn = 1.5$, $\cfp = 1.0$, $N = 1{,}000$). Our bound is uniformly tighter, with increasing advantage at larger $\rhostar$.}
    \label{tab:bound}
    \begin{center}
      \begin{small}
        \begin{tabular}{ccccc}
          \toprule
          $\rhostar$ & Risk  & $\Bpw$ bound & $\Btv$ bound & Improv. \\
          \midrule
          2.0        & 0.243 & 0.514        & 0.701        & 27\%    \\
          5.0        & 0.161 & 0.870        & 1.368        & 36\%    \\
          10.0       & 0.105 & 1.124        & 1.844        & 39\%    \\
          \bottomrule
        \end{tabular}
      \end{small}
    \end{center}
  \end{table}

  \begin{figure}[h]
    \centering
    \includegraphics[width=0.7\textwidth]{figures/fig1_bound_tightness.png}
    \caption{Bound comparison across shift ratios (synthetic). The prevalence-weighted bound ($\Bpw$, orange) is 27--39\% tighter than the TV bound ($\Btv$, grey) while both remain conservative relative to realised risk (blue).}
    \label{fig:bound}
  \end{figure}

  \subsection{Bayesian Interval Coverage}

  To evaluate robustness to BBSE estimation error, we measure the frequentist coverage of the $95\%$ Bayesian credible interval across $\rhostar \in \{1, 1.5, 2, 3, 5\}$ (50 seeds each).
  Coverage is $98\%$ at $\rhostar = 1.0$, $96\%$ at $\rhostar = 1.5$, $90$--$92\%$ at $\rhostar = 2$--$3$, and $84\%$ at $\rhostar = 5.0$ (extreme shift), with near-zero bias ($\leq +0.03$).
  The interval widens appropriately as shift grows ($0.22$ at $\rhostar = 1$ to $0.88$ at $\rhostar = 5$).

  \subsection{Safety--Utility Tradeoff}

  \begin{figure}[h]
    \centering
    \includegraphics[width=\textwidth]{figures/fig3_recal_efficiency.png}
    \caption{Label-free recalibration safety--utility tradeoff (synthetic). \textbf{Left:} wider uncertainty interval $\delta$ increases the negative prediction rate (more conservative). \textbf{Right:} realised risk stays well below $\alpha = 0.3$ across assumed shift ratios. The guarantee holds at every point.}
    \label{fig:efficiency}
  \end{figure}

  \begin{table}[h]
    \caption{Label-free recalibration safety--utility tradeoff as interval half-width $\delta$ varies ($\rhohat = 1.0$, $\alpha = 0.30$, synthetic data). \textbf{Take-away:} the tradeoff is smooth; the guarantee holds at every point.}
    \label{tab:tradeoff}
    \begin{center}
      \begin{small}
        \begin{tabular}{cccc}
          \toprule
          Half-width $\delta$ & $\asafe$ & Realised risk & Positive rate \\
          \midrule
          0.05                & 0.286    & 0.200         & 0.372         \\
          0.10                & 0.273    & 0.191         & 0.381         \\
          0.20                & 0.246    & 0.172         & 0.402         \\
          0.30                & 0.217    & 0.150         & 0.426         \\
          0.50                & 0.157    & 0.112         & 0.466         \\
          0.80                & 0.123    & 0.092         & 0.491         \\
          \bottomrule
        \end{tabular}
      \end{small}
    \end{center}
  \end{table}

  \subsection{Cost Sensitivity}

  \begin{figure}[h]
    \centering
    \includegraphics[width=0.5\textwidth]{figures/fig4_cost_sensitivity.png}
    \caption{Effect of cost ratio $\cfn / \cfp$ on $\Bpw$ and $\asafe$ (synthetic). \textbf{Take-away:} as cost asymmetry grows, $\Bpw$ increases linearly and $\asafe$ shrinks; recalibration becomes infeasible at extreme ratios.}
    \label{fig:cost}
  \end{figure}

  \begin{table}[h]
    \caption{Effect of cost ratio $\cfn / \cfp$ on the prevalence-weighted bound and $\asafe$ ($\pii{1}{tr} = 0.20$, $\delta = 0.2$, $N = 1{,}000$). \textbf{Take-away:} recalibration remains feasible up to $\cfn / \cfp = 20$; beyond that, $\asafe$ approaches zero.}
    \label{tab:cost}
    \begin{center}
      \begin{small}
        \begin{tabular}{cccc}
          \toprule
          $\cfn / \cfp$ & $\Bpw$ bound & $\asafe$ & Feasible? \\
          \midrule
          2             & 0.870        & 0.246    & Yes       \\
          5             & 1.012        & 0.232    & Yes       \\
          10            & 1.724        & 0.165    & Yes       \\
          15            & 2.436        & 0.097    & Yes       \\
          20            & 3.148        & 0.029    & Yes       \\
          \bottomrule
        \end{tabular}
      \end{small}
    \end{center}
  \end{table}

  \subsection{Exchangeability Sanity Check}

  \begin{figure}[h]
    \centering
    \includegraphics[width=0.7\textwidth]{figures/fig_exchangeability.png}
    \caption{Cost-weighted risk under exchangeability ($\rhostar = 1.0$, Folktables, 30 seeds). \textbf{Take-away:} Ours (label-free) has the largest safety margin (risk $0.244$ vs.\ $\alpha = 0.3$); other methods operate near the boundary with risk $\approx 0.30$.}
    \label{fig:exchangeability}
  \end{figure}

  \subsection{Method Breakdown Under Increasing Shift}

  \begin{figure}[h]
    \centering
    \includegraphics[width=0.5\textwidth]{figures/fig2_breakdown_curve.png}
    \caption{Realised cost-weighted risk as $\rhostar$ increases (synthetic). \textbf{Take-away:} Standard and Symmetric CRC risk drops (cost asymmetry effect); Oracle/BBSE maintain risk near $\alpha$; ours (label-free) decreases most aggressively due to conservative calibration.}
    \label{fig:breakdown}
  \end{figure}

  \section{Reproducibility}
  \label{app:reproducibility}

  \textbf{Code.}
  All experiments, tests, and figure generation are contained in a single repository.
  The master runner \texttt{run\_all.py} executes the full pipeline: unit tests $\to$ PACPS replication $\to$ Folktables ACSIncome (exp0, exp6) $\to$ ACSPublicCoverage (exp8) $\to$ deferral curve (exp7) $\to$ synthetic experiments (exp1--exp4) $\to$ GAIA agentic experiment (exp12) $\to$ figure and table generation.
  A single command reproduces all results:
  \begin{verbatim}
  uv run python run_all.py
\end{verbatim}

  \textbf{Data.}
  Folktables uses the 2018 American Community Survey via the \texttt{folktables} Python package, which downloads data automatically.
  The PACPS Heart Disease benchmark uses the CDC Heart archive following the protocol of \citet{fisch2022conformal}.

  \textbf{Dependencies.}
  Python $\geq 3.13$, managed via \texttt{uv}.
  Core: NumPy, SciPy, pandas, scikit-learn, matplotlib, seaborn.
  GAIA: OpenAI API (GPT-4o-mini for agent and judge), HuggingFace \texttt{datasets}.
  All dependencies are pinned in \texttt{uv.lock}.

  \textbf{Computational requirements.}
  Folktables experiments and all synthetic experiments complete in under 5 minutes on a single CPU core.

  \textbf{Random seeds.}
  All experiments use fixed seeds (0--29 for 30-seed experiments).
  The Bayesian shift estimation uses seed-controlled Monte Carlo sampling (10{,}000 draws).

  \textbf{Test suite.}
  58 unit tests cover: CRC validity under exchangeability, bound non-negativity, BBSE consistency, recalibration sanity, Bayesian posterior correctness, credible interval coverage ($\geq 85\%$ over 100 simulated trials), three-way decision zone properties, deferral rate arithmetic, GAIA episode schema, judge score parsing, and policy action mapping.

\end{document}
